% Folder 142, batch 06 (archivists' pages 101-120).
% Pass: Opus 5.5 (claude-opus-5-5), 2026-10-02 — first pass, unchecked against the pages by a human.
%
\documentclass[11pt,a4paper]{article}
\input{../preamble/grothendieck}

\folder{142}
\batch{6}
\pages{101}{120}
\foldertitle{Action de Galois sur Teichmüller : notes manuscrites (s.d.)}
\dating{[à partir de 1978]}
\watermark{Édition de démonstration}

\begin{document}

\page{101}
\note{la page continue une discussion commencée avant ce lot : les conditions 2°, 3°, 4° auxquelles elle renvoie sont énoncées dans les pages précédentes.}
À vrai dire, il ne serait pas difficile d'\uncertain{expliciter} une condition 5°) qui mette ensemble 3°) et 4°), quand on se donne \ill{} $X'_1 \to X'_2$, un \uncertain{morphisme}
\[
  X'_1 \xrightarrow{\;f_1\;} Y, \qquad X'_2 \xrightarrow{\;f_2\;} Y
\]
\note{au-dessus de $X'_1 \to X'_2$, une flèche biffée et surchargée, et une flèche oblique $f_1$ vers $Y$ ; le petit diagramme est trop raturé pour être rendu fidèlement.}
\ill{} $f_2$ \struck{\ill{}} \uncertain{induisant} \uncertain{des homomorphismes} \struck{\uncertain{fidèles}} sur les $\pi_1$, de sorte que \uncertain{surjectif} sur les \uncertain{fibres} \uncertain{pour} \uncertain{l'application} \ill{}
$\hat\pi_1(X'_2, I'_2)$, de \uncertain{partition} \ill{} \struck{\ill{}} $\Theta'_2$ de $\Theta$ sur $\hat\pi_1(Y, f_2(I'_2))$, \uncertain{soit} que l'\uncertain{action} \ill{} \uncertain{respecte} les dessins \uncertain{pour} : $\hat\pi_1(Y, f_2(I'_2))$ \ldots
Mais je ne \uncertain{sais} pas \uncertain{encore} si j'aurai besoin de cette généralisation \uncertain{commune} de 3° et 4° (cet $=$ de 2°, 3°, 4°, à \uncertain{mon} \uncertain{avis} !).

Les conditions 2° : 4° \uncertain{sont} \uncertain{pour} \uncertain{moi} \uncertain{seulement} \uncertain{nécessaires} pour \uncertain{trouver} \ill{} \uncertain{des} conditions \uncertain{nécessaires} sur les \uncertain{paramètres} \uncertain{introduits} qui \uncertain{répondent} \uncertain{aux} \ill{} pour que ils \uncertain{soient} \uncertain{effectifs}, i.e. \uncertain{proviennent} tel \ill{} et \ill{} d'un \ill{}. \uncertain{Le} \uncertain{problème} \uncertain{de} \uncertain{trouver} \uncertain{l'ensemble} \uncertain{naturellement} \ill{} de \uncertain{dégager} \uncertain{des} conditions nécessaires et \uncertain{suffisantes} qui \uncertain{permettent} de \uncertain{caractériser} les $\Theta$ provenant de $\Gamma$ \ldots

\page{102}
Soit $X/K$, et soit $T$ une partie de $X(\mathbb{C})$ \add{$X =$ \ill{}}, loc\supplied{alemen}t fermée 1-connexe. Soit $X' \to X$ \struck{un \ill{}} \add{un} \struck{\uncertain{voisinage}} \ill{} $K$-\uncertain{multiplicité}, qui sur \struck{les \uncertain{points} de} un voisinage de $T$ soit fini étale. Alors $X' \mid T = X'_T$ \ill{} un ens\supplied{emble} fini de copies de $T$ ; il y a une relation d'équivalence des \ill{} \ill{} la $T$-adjacence de points de $X'$ au-dessus de $T$. On se donne \struck{$X'_T$} $x', y' \in X'_T$ $T$-adjacents ;
\marginal{on \uncertain{note} \uncertain{$W'$} \ldots{} $x', y' \in X'_{\bar K}$}
si $u \in \Gamma$, on a $u(x')$, $u(y')$ \struck{\ill{}} sont aussi $T$-adjacents. Mais si \uncertain{on} les \ill{}, on \uncertain{trouve} \struck{$\hat\pi_1(X'_{\bar K}, \ldots)$} \add{\uncertain{dans} $T$-horizontal} \ill{}. \uncertain{On} \uncertain{peut} regarder $l^{T}_{x';y'} : x' \to y'$ dans \ill{} et les \uncertain{transformés}, \struck{\ill{}} \ill{} :
\struck{$u(l^{T}_{y';x'})$} et la \struck{\uncertain{composition}} : $l^{T}_{u(y'),u(x')}$ \uncertain{qui} \uncertain{lie} $u(x'), u(y')$ (d'homotopie \uncertain{horizontale}) défini par \uncertain{lui}.
\begin{equation*}
  \tag{53} u\, l^{T}_{y';x'} \;=\; g^{T}_{x,y}(u)\; l^{T}_{u(y'),u(x')}
\end{equation*}
où
\[
  g^{T}_{x',y'}(u) \in \hat\pi_1\bigl(X',\, T(u(x'),u(y'))\bigr)
\]
\marginal{$T(u(x'),u(y'))$ : structure horizontale de $X'_T$ qui contient $u(x')$, $u(y')$}
\note{l'indice de $g$ dans (53) se lit $x,y$ ; dans la ligne suivante il porte un exposant raturé et se lit $x',y'$.}

Ceci posé, on \uncertain{voit}
\begin{equation*}
  \tag{54} \hat f\bigl(\underbrace{g^{T}_{y';x'}(u)}_{\in\, \hat\pi_1(X,T)}\bigr) \;=\; g^{T}_{y,x}(u)
\end{equation*}

\page{103}
\uncertain{Prenons} que \struck{\ill{}} \add{\uncertain{les composantes connexes}} $T_\alpha$ de $X'_T$, considérons
\[
  \hat\varphi_\alpha : \hat\pi_1(X'_{\bar K}, T_\alpha) \longrightarrow \hat\pi_1(X, T)
\]
(c'est un hom\supplied{omorphisme} injectif si $X'$ rev\supplied{êtement} étale de $X$), et tous les \struck{\uncertain{éléments} \ill{}} images des $\varphi_\alpha$ sont les conjugués \add{si \struck{\ill{}} \uncertain{l'on} \uncertain{prend} des \uncertain{points} $x'_\alpha, y'_\alpha$ \uncertain{convenables}} des \ill{} de l'\ill{} \ill{}. Donc (la formule \uncertain{(54)} \add{\uncertain{pour} \uncertain{tout} $\alpha$}) \ill{} \ill{} :
\[
  g^{T}_{y,x}(u) \in \bigcap_{\alpha \in \pi_0(X'_T)} \operatorname{Im} \hat\varphi_\alpha
  \qquad \forall u \in \Gamma = \Gamma_{\bar K/K}
\]
une condition \uncertain{utile} \uncertain{nécessaire}, \uncertain{issue} \uncertain{de} la \uncertain{variété} (\uncertain{finie} \uncertain{topologique} \ill{}) $X'$ \uncertain{puisque} \struck{\uncertain{particulière}} \uncertain{étale} \ill{} d'une \ill{} $K$-\uncertain{schéma}, pour laquelle $f : X' \to X$ soit \uncertain{net} $/K$, et telle que $\forall \alpha \in \pi_0(X'_T)$, \struck{$x'_\alpha, y'_\alpha$} $\forall u \in \Gamma$, $u(x'_\alpha)$ et $u(y'_\alpha)$ soient $T$-adjacents (ce qui est une condition nettement plus faible que celle que $x', y' \in X'_{\bar K}$ \add{i.e.\ $x', y'$} \uncertain{soient} fixés sous $\Gamma$)\ldots

\marginal{\uncertain{Si} $X'$ \uncertain{est} \uncertain{un} \uncertain{rev.} de $X$, l'ens\supplied{emble} $E = \pi_0(X'_T)$ des \uncertain{« feuilles »} \uncertain{(55)} qui \ill{} de $X'$ \ill{}, $\pi = \pi_1(X,T)$, \uncertain{est} la \uncertain{« fibre »} \uncertain{définie} par $T \to X$ ($T$ 1-connexe), \ill{} une action de $\pi$ sur $E$, et \ill{} \uncertain{dist.} \ill{} $\pi \to \mathfrak{S}_E$ \ill{} \ill{}}
\note{note marginale écrite en diagonale dans l'angle inférieur gauche de la page ; lecture très partielle.}

\section{Groupoïde fondamental d'une carte trouée (préliminaires à l'action de Galois)}
\note{titre écrit de sa main sur une feuille de garde, p.~104, avec la mention « (14 pages) » ; la feuille ne porte rien d'autre.}

\page{105}
\note{page de quatre figures : la sphère privée de $0, 1, \infty$, chaque point entouré d'un petit cercle, avec des sommets marqués sur ces cercles et des chemins orientés entre eux ; ce sont des présentations du groupoïde fondamental relatif à divers ensembles de points de base.}

\drawing{Première figure. Trois petits cercles autour de $0$, $1$, $\infty$, deux sommets sur chacun ; les arcs des cercles portent $\lambda^0_+, \lambda^0_-$, $\lambda^1_+, \lambda^1_-$, $\lambda^\infty_+, \lambda^\infty_-$ ; trois chemins $a^1$ (de $\infty$ vers $0$), $a^\infty$ (de $0$ vers $1$), $a^0$ (de $1$ vers $\infty$). Légende : « 6 sommets, 9 flèches, \struck{1} 2 relations ».}

\drawing{Deuxième figure. Même sphère ; un seul sommet sur chacun des cercles de $0$ et $\infty$, deux sur celui de $1$ ; flèches $\Lambda^\infty_+$, $\Lambda^0_+$, $\lambda^1_+$, $\lambda^1_-$, $a^\infty$ (de $0$ vers $1$), $a^0$ (de $1$ vers $\infty$). Légende : « 4 sommets, 1 relation, 6 flèches » ; dessous : « 3 cas ».}

\drawing{Troisième figure, biffée de traits en zigzag : cercles autour de $\infty$, $0$, $1$, un chemin de $1$ vers $\infty$ et un de $0$ vers $1$ ; au-dessus, « 1 cas ».}

\drawing{Quatrième figure. Cercles autour de $0$ et $1$, un chemin $a^\infty$ de $0$ vers $1$, le point $\infty$ isolé dans un petit cercle. Inscriptions $\Lambda^0_+ = (\Lambda^0_-)^{-1}$, $\Lambda'^1_- = (\Lambda'^1_+)^{-1}$. Légende : « 2 sommets, trois flèches, 0 relation » ; dessous : « 3 cas ».}

\drawing{Cinquième figure. Un seul sommet sur le cercle de $0$, un chemin vers le cercle de $1$ (étiquette raturée), le point $\infty$ isolé ; flèches $\Lambda^+_0$ et $\Lambda^{01}_+$. Légende : « 1 sommet, 2 flèches, 0 relation » ; dessous : « 3 cas. »}

\page{107}
\drawing{Au crayon, deux sphères. La première : pôles $0$ (en haut) et $\infty$ (en bas), méridiens et équateur ; sur l'équateur, quatre petits cercles orientés autour de $1$, $i$, $-1$, $-i$ (étiquettes partiellement lisibles), reliés aux pôles par des chemins fléchés, certains en pointillé.}

À droite de la première sphère :
\[
  z \mapsto -z, \qquad z \mapsto \frac{1}{z}, \qquad \text{\struck{\ill{}}}
\]
\[
  1 \longrightarrow \mathbb{F}_2^{\,2} \longrightarrow \mathfrak{A}_4 \longrightarrow \mathbb{Z}/3\mathbb{Z} \longrightarrow 1
\]

\drawing{La seconde sphère : petit cercle orienté autour du pôle supérieur et autour de $\infty$ en bas ; sur l'équateur, cercles autour de $1$ et $-1$, points $-i$ et $i$ ; chemins marqués $\beta$, « $\beta$ relevé », « $\gamma$ relevé » (deux fois), le dernier avec la mention « $\gamma$ relevé, symétrique de \ill{} ».}

\[
  l_\infty\, l_1\, l_0\, l_{-1} = 1
\]
\note{dans le symbole $\mathfrak{A}_4$ la lettre est une capitale cursive, lue comme le groupe alterné.}

\page{109}
\note{la feuille porte deux couches d'écriture tête-bêche : en haut, à l'endroit, des formules sur les sous-groupes ; le reste de la page, retourné, est occupé par des dessins de sphères. On donne d'abord les formules, puis les dessins.}

\[
  E_0 \;\text{---}\; \pi_0(E_0) \;\text{\struck{$\simeq E_0/\pi \;\mathrm{Ssgr}($}}
\]

\begin{tikzcd}
  E_0 \arrow[r, no head] \arrow[d] & \pi_0(E_0) \arrow[d] \\
  \mathrm{Ssgr}(\pi) \arrow[r] & \mathrm{Ssgr}(E)/\pi
\end{tikzcd}

\note{les deux flèches verticales descendent de $E_0$ et de $\pi_0(E_0)$ ; sous la ligne du bas, une accolade porte l'inscription « $\Sigma$ opérant ci-dessus » (lecture incertaine, la lettre surchargée).}

\[
  \Sigma \;(= \Gamma \cdot \pi) \supset \pi \;\text{---}\; E_0 \longrightarrow \mathrm{Ssgr}(\pi)/\pi,
  \qquad \Sigma/\pi = \Gamma
\]
\note{une flèche descend de $\Sigma$ vers $\mathrm{Ssgr}(\pi)/\pi$.}

\drawing{Partie retournée. Une sphère avec pôles $0$ et $\infty$, l'équateur passant par $1$, quatre petits cercles sur l'équateur ($a$, $b$ marqués) ; deux flèches en partent vers le bas, avec la légende « division par $D$ (multiplication par les \ill{} \uncertain{$R_s$}) ». Une seconde flèche « division par $D^+$ (\uncertain{multiplication} \ill{} \uncertain{sur} $(0, \infty)$) » mène à une sphère dont l'équateur passe par $1$.}

\drawing{Sphère centrale, au crayon et à l'encre : sur l'équateur le point $Q_1$ près de $-1$, le cercle autour de $1$ avec $s^+_1$, $s^-_1$, $a^{++}_1$, $a^{+-}_1$ ; chemins $b^+_1$, $b^-_1$ vers les cercles des pôles ($R^-_\infty$, $a^-_\infty$, $R^+_\infty$, $a^+_\infty$) ; inscriptions $d^+_1 (c^+_1)^{-1}$, $(a^-_1 a^+_1)^{-1}$, $d^-_1 (c^-_1)^{-1}$ ; un axe traverse la sphère, une flèche de rotation à côté.}

\drawing{Sphère inférieure : pôles marqués $-\zeta$ et $-\bar\zeta$ (lecture incertaine), cercles autour de $\infty$ (point $R^+_\infty$), de $0$ (point $R_0$) et de $1$ ; chemins $c^\pm_1$, $d^\pm_1$, $b^\pm_1$, $a^{\pm\pm}_1$, $s^\pm_1$, le point $Q_1$. À côté : « \uncertain{stable} par \uncertain{symétrie} \uncertain{autour} de l'axe ».}

\drawing{À gauche, deux sphères raturées ; sur la seconde, les points $0$, $1$, $\infty$ sur un équateur, une moitié hachurée. À côté, la table
$0, \infty \mapsto \infty$, $1 \mapsto 0$, $-1 \mapsto 1$.}

\page{110}
\[
  \Sigma_J = \mathbb{P}\bigl(V(J)\bigr) \simeq \check{\mathbb{P}}\bigl(V(J)\bigr)
\]
\note{le premier $\mathbb{P}$ est récrit sur une autre lettre.}

\struck{Suite} \quad $0 \to F \to V(J)_\Sigma \to \mathcal{O}_\Sigma(1) \to 0$
\qquad
$0 \to V(J) \to \mathcal{O}^J \xrightarrow{\text{somme}} \mathcal{O} \to 0$

\[
  F(1) \simeq \mathcal{O}(\omega)
\]
($\mathcal{O}$ tordu par $\underline{\omega} = \underline{\omega}(J)$, i.e.\ $\det V(J)$)
\[
  \check{\mathcal{O}}(1) \simeq \check F(\omega)
\]
(par $F \wedge \mathcal{O}_\Sigma(1) \xrightarrow{\sim} \det V(J)_\Sigma \simeq \mathcal{O}_\Sigma(\omega_\Sigma)$, prod.\ ext.) \marginal{attention signe !}
\[
  \Omega^1_\Sigma \simeq F(-1) \simeq \mathcal{O}(-2)(\omega), \qquad
  \check\Omega^1_\Sigma \simeq \mathcal{O}(2)(\omega) \quad \text{fibré tangent à } \Sigma
\]

\marginal{\uncertain{Accouplement} $F \times V(J) \to \underline{\mathcal{O}}$ \ill{} donné par \ill{} \uncertain{(isomorphisme)} $x \wedge y \wedge \ldots = \lambda \, x \wedge y \wedge \ldots$, $x \in F$, $y \in V(J)$ \ill{} i.e.\ \ill{} \ill{} de $(x, y)$.}
\note{note marginale en diagonale, à gauche du haut de la page ; lecture très partielle.}

Points $R_i$ ($i \in J$) (« trous ») de $\Sigma$ et $\mathbb{P}\Sigma$ ; $\tilde R_i$ \add{\uncertain{(entouré)}}
$R_i \subset$ \uncertain{$\mathbb{V}(J)$} \qquad $\mathcal{O}(1)_{R_i} \simeq \mathcal{O}$
$\tilde R_i$ pt au-dessus de $R_i$, image de $e_i - e_j = \varphi_{ji}$ ($\equiv e_i - e_k$) dans $V(J)/R_i$ ;
$R_i = \mathcal{O}_k(e_k - e_j) \simeq \mathcal{O}(\underline\omega)$ \note{l'indice $k$ de $\mathcal{O}$ est douteux.}
\uncertain{on a choisi} $e_k - e_j$, l'orientation $(i, j, k)$.
\struck{Désignons par $\tilde R_i$ \ldots{} dans $\mathbb{P}\Sigma_J = \check{\mathbb{V}}^*(\mathcal{O}_\Sigma(1))$ l'élément $1$ au-dessus de $R_i$.}
\[
  \Gamma\bigl(\Sigma_J, \mathcal{O}_\Sigma(1)\bigr) \simeq V(J)
  \qquad \text{\struck{$e_j - e_k$ section $\omega_{j,k}$}}
\]
\[
  \varphi_{jk} = e_k - e_j \in \Gamma\bigl(\Sigma, \mathcal{O}_\Sigma(1)\bigr) \quad \text{nulle en } R_i,
  \qquad \varphi_{kj} = -\varphi_{jk}
\]
\marginal{$\varphi_{ij} + \varphi_{jk} + \varphi_{ki} = 0$}
\[
  \varphi_{jk}(R_i) = 0 \qquad
  \boxed{\varphi_{jk}(R_k) = \tilde R_k, \quad \varphi_{jk}(R_j) = -\tilde R_j}
\]
\marginal{attention signes !}

Points $Q_i$ ($i \in J$), $Q^\omega_i$ ($i \in J$, $\omega \in \underline\omega$)
$Q_i \subset$ \uncertain{$\mathbb{V}(J)$} $= \mathcal{O}\bigl(-2e_i + (e_j + e_k)\bigr)$
\qquad $\mathcal{O}(1)_{Q_i} \simeq \mathcal{O}(\underline\omega)$
les deux points privilégiés de $\mathcal{O}(\underline\omega)^*_{Q_i}$, indexés par $\underline\omega$, sont notés $\tilde Q_i, \tilde Q'_i$, ou $Q_{i,\omega}$ ($(i,\omega) \in J \times \underline\omega$)

($e_i - e_j = e_{\omega i} - e_{\omega^2 i}$ dans $V(J)$ \uncertain{mod} $Q_i$)
donc \ill{} \struck{pour $\varphi_{\omega i, \omega' i}(Q_i) =$}
\[
  \tilde Q_i = \varphi_{jk}(Q_i), \qquad \tilde Q'_i = -\varphi_{jk}(Q_i) = \varphi_{kj}(Q_i)
\]
\[
  \boxed{Q^\omega_{i} = \varphi_{\omega i, \omega' i}(Q_i)}
\]
\note{la formule encadrée est récrite plusieurs fois ; l'indice de $\varphi$ est en partie noyé dans l'encre.}
\marginal{encadré : $Q^{\omega+1}_i = -Q^\omega_i$ ; $\varphi_{\omega i, \omega' i}(Q_i) = \ldots$ \ill{}}

Points $P_\omega$ ($\omega \in \underline\omega$), $P^\omega_i$ ($i \in J$, $\omega \in \underline\omega$). Soit $\zeta$ une racine primitive cubique de $1$ (on prendra $\zeta = \exp 2i\pi/3$).

Alors

\page{111}
\note{les ensembles $\Sigma$ portent en exposant de petites piles de signes ($\ast$ et points) qui disent quels points ont été ôtés ; on les rend par $\ast$ et $\cdot$ dans l'ordre où ils se lisent, sans garantir la disposition.}

Carte régulière (\uncertain{orientée}) des polygones réguliers $P_n$ ($n \geq 2$) en la trouant \add{\uncertain{(au-dessus de)}} les sommets et les centres des faces \emph{et} les \uncertain{milieux} \add{\uncertain{\ill{}}} des segments
\note{le « et » est souligné par lui.}
\[
  \Sigma_n^{\ast\ast} = \Sigma_n \setminus \bigl\{ S(\Sigma_n) \cup A(\Sigma_n) \cup F(\Sigma_n) \bigr\}
\]
\marginal{\uncertain{centres faces}}
\uncertain{divisions} par $D_n^+ \subset D_{2n}^+$
\[
  \Sigma_n^{\ast\ast} \xrightarrow{\ \deg n\ } \Sigma_1^{\ast\ast} = \mathbb{P}_1 \setminus \{\underbrace{0, \infty}_{\text{« pôles »}}, 1, -1\}
  \xrightarrow{\ \deg 2\ } \Sigma_3^{\ast\cdot} = U_{0,3}
\]
\marginal{division par $D_n/D_n^+$}
\marginal{(\uncertain{cartes} régulières $\Sigma_\alpha$ \uncertain{à} \uncertain{sommets})}
\[
  \Sigma_n^{\ast\ast} \simeq \Sigma_1^{\ast\ast}/D_n \;\; \text{\uncertain{(sic)}}
\]
\note{sous $\Sigma_n^{\ast\ast}$, un signe $\simeq$ vertical renvoie à $\Sigma_{2n}^{\ast\cdot}$ ; à droite, $\Sigma_n^{\ast\ast}/D_n$ est relié à la ligne précédente par $\parallel$. La disposition exacte est incertaine.}
\[
  \Sigma_{2n}^{\ast\cdot} \longrightarrow \Sigma_1^{\ast\cdot} \xrightarrow{\ 2\ } \Sigma_1^{\ast\cdot} \simeq U_0^3
\]
\marginal{division par $D_{2n}^+$}

\marginal{(\uncertain{on n'a} \uncertain{troué} que les sommets et centres faces) À cause de \uncertain{l'isomorphisme} $D_{2n}$ \uncertain{opérant} sur $\Sigma_n^{\ast\cdot}$ ($\simeq \Sigma_{2n}^{\ast\cdot}$), on \uncertain{peut} \ill{}}

Diagramme de quotients :

\begin{tikzcd}[column sep=small, row sep=small, nodes={font=\scriptsize}]
  & \Sigma_n^{\ast\cdot} \simeq \Sigma_{2n}^{\ast\cdot} \arrow[dl] \arrow[dr] & \\
  \Sigma_{2n}^{\ast\cdot}/D_{2n}^+ \arrow[d] & & \Sigma_n^{\ast\cdot}/D_n^+ \\
  \Sigma_{2n}^{\ast\cdot}/D_{2n} & &
\end{tikzcd}

\note{le diagramme principal est à droite, entouré d'un trait :}

\begin{tikzcd}[column sep=small, row sep=small, nodes={font=\scriptsize}]
  & \Sigma_n^{\ast\cdot} \simeq \Sigma_{2n}^{\ast\cdot} \arrow[d] & \\
  & \Sigma_n^{\ast\cdot}/D_n^+ \simeq \Sigma_1^{\ast\ast} \simeq \Sigma_2^{\ast\cdot} \arrow[dl] \arrow[dr] & \\
  \Sigma_n^{\ast\ast}/D_{2n}^+ \arrow[dr] & & \Sigma_n^{\ast\cdot}/D_n \arrow[dl] \\
  & \Sigma_n^{\ast\cdot}/D_{2n} &
\end{tikzcd}

\note{au-dessus du diagramme entouré, $\Sigma^{\ast\ast}/D_n$, relié par une flèche venant de la gauche.}
\[
  \Sigma_n^{\ast\ast} \xrightarrow{\ D_n^+\ } \Sigma_1^{\ast\cdot} \simeq U_0^3
\]
\uncertain{si} \uncertain{on} \uncertain{divise} par $D_n$ \uncertain{tout} \uncertain{entier}, on \uncertain{trouve} \ill{} \uncertain{à} \uncertain{une} \uncertain{réflexion} \uncertain{près} \uncertain{les} \uncertain{arêtes}.

\page{112}
\note{la page reprend les notations de la p.~110 : points $R_i$, $Q_i$, $Q^\omega_i$ ; elle est faite surtout de figures.}

\drawing{Au crayon, trois bandes concentriques qui figurent le revêtement au-dessus de $\Sigma$ ; on y lit les points $Q_1$, $Q_2$, $Q_3$ (marqués d'une croix), $Q^{\omega'}_1$, $Q^{\omega^2}_1$, $Q^{\omega'}_3$, $Q^{\omega}_3$, et, autour de chaque trou $R_i$, les points $R_{i,\omega}$, $R_{i,\omega'}$, $\tilde R_i$, $-\tilde R_i$, $R^\omega_i$, $R^{\omega'}_i$, $-R^\omega_i$, $-R^{\omega'}_i$. Autour de $R_1$ et de $R_2$, à l'encre, de petits quadrilatères de flèches courbes reliant ces points.}

\drawing{Un cercle, le centre relié par des rayons au crayon à douze points du bord, nommés dans l'ordre
$R^{\omega'}_2$, $-R^\omega_1$, $-R^{\omega'}_1$, $R^\omega_3$, $R^{\omega'}_3$, $-R^\omega_2$, $-R^{\omega'}_2$, $R^\omega_1$, $R^{\omega'}_1$, $-R^\omega_3$, $-R^{\omega'}_3$, $R^\omega_2$ ;
des flèches courbes relient les points voisins deux à deux.}

\drawing{Un petit cercle avec les points $1$, $-1$ et $\zeta$, l'arc de $1$ à $\zeta$ fléché.}

\[
  Q^\omega_i = R_j - R_k \longmapsto R_j - \zeta R_k = P^\omega_i \;\text{\struck{\ill{}}}\; = P^{\omega,\zeta}_i
\]
\[
  P^{\omega,\zeta}_i = R_{\omega i} \;\text{\struck{$\zeta$}}\; - \zeta R_{\omega^2 i} \;\equiv\; -P^{\omega,\zeta}_{\omega^2 i}
\]
\[
  P^{\omega',\zeta'}_i = R_{\omega' i} - \zeta' R_{\omega i} = R_k - \zeta^2 R_j = -\zeta' P^{\omega,\zeta}_i
\]
\note{les exposants $\omega^4 i$, $\omega^2 i$ de la deuxième ligne sont en partie surchargés ; lecture incertaine.}

$P^{\omega,\zeta}_i$ prend 12 valeurs distinctes ;
\uncertain{$P_{i}$} six valeurs pour \struck{\ill{}} $\zeta$ fixé (les $P^\omega_i$) et les six \ill{}

\page{113}
\note{p.~1 de l'auteur : le chiffre est en tête de la page ; la page suivante porte « 2 ».}

Carte $\mathcal{X}$ (pas orientée) \add{\uncertain{finie}}, $\mathcal{R}$ ens\supplied{emble}. des repères, $S = S_0 \sqcup S_1 \sqcup S_\infty$ \add{sommets, arêtes, faces}
l'ens\supplied{emble} de ces feuillets, identifiés \struck{aux} à leurs centres, dans \ill{} \ill{} désignons par $\mathcal{X}$ aussi le \struck{\ill{}} \uncertain{complexe} topologique correspondant : $\mathcal{X}$,
\[
  S = S_0 \sqcup S_1 \sqcup S_\infty \subset \mathcal{X}
\]
si $\varphi \in S$, on désigne par $R_\varphi$ le sommet associé. On identifie un $r = (s, a, f) \in \mathcal{R}$ à la région triangulaire correspondante (c'est un triangle de la subdivision barycentrique), de sommets $R_s$, $R_a$, $R_f$, d'indices respectifs $0, 1, \infty$. On oriente un triangle dans le sens $0, 1, \infty$ (orientation du repère). Les repères « adjacents » à $r$
\[
  \tau_0(r),\ \tau_1(r),\ \tau_\infty(r),
\]
opposés respectivement aux sommets d'indice $0$, $1$, $\infty$ de $r$. Notons
\begin{equation*}
  \tag{1}
  \begin{cases}
    \tau_1(\tau_0(r)) = \rho_\infty(r) \\
    \tau_0(\tau_1(r)) = \rho_\infty^{-1}(r)
  \end{cases}
  \qquad \text{(face conservée)}
\end{equation*}
déduits de $r$ par « rotation » \add{d'un cran} autour \uncertain{des} \uncertain{sommets}, \uncertain{du} type $\infty$ vers $0$, dans le sens \struck{indiqué} de l'orientation définie par $r$ resp\supplied{ectivement}. dans le sens \uncertain{inverse}. De même
\begin{equation*}
  \tag{2}
  \begin{cases}
    \tau_\infty \tau_1(r) = \rho_0(r) \\
    \tau_1 \tau_\infty(r) = \rho_0^{-1}(r)
  \end{cases}
  \qquad \text{(sommet conservé)}
\end{equation*}
où $\rho_0$ est la rotation \add{d'un cran} autour du sommet de type $0$, dans le sens indiqué par l'orientation définie par $r$. Enfin
\begin{equation*}
  \tag{3}
  \tau_0 \tau_\infty(r) = \tau_\infty \tau_0(r) \overset{\text{def}}{=} \rho_1(r)
\end{equation*}

\drawing{À gauche, un morceau de carte et sa subdivision barycentrique : sommets marqués $0$, $1$, $\infty$, le triangle central $r$ (entouré) avec ses sommets $R^0_r$, $R^1_r$, $R^\infty_r$ et ses côtés $a^0_r$, $a^1_r$, $a^\infty_r$, $\lambda^0_r$, $\lambda^1_r$, $\lambda^\infty_r$ ; les triangles voisins sont étiquetés $\tau_0(r)$, $\tau_1(r)$, $\tau_\infty(r)$ (signe $-$), $\rho_\infty(r) = \tau_1\tau_0(r)$, $\rho_\infty^{-1}(r) = \tau_0(\tau_1(r))$, $\tau_\infty\tau_1(r) = \rho_0(r)$, $\tau_1\tau_\infty(r) = \rho_0^{-1}(r)$, $\rho_1(r) = \tau_0\tau_\infty(r) = \tau_\infty\tau_0(r)$ (signe $+$) ; les points $R_f$ ($\infty$), $R_s$ ($0$), $R_a$ ($1$), $R^{0,1}_r$, $R^{1,\infty}_r$, $R^{0,\infty}_r$ ; flèches courbes de rotation autour des sommets.}

[Pour mémoire, les relations dans le groupe cartographique
\begin{equation*}
  \tag{4}
  \tau_0^2 = \tau_1^2 = \tau_\infty^2 = (\tau_0 \tau_\infty)^2 = 1
\end{equation*}
\begin{equation*}
  \tag{5}
  \begin{cases}
    \rho_\infty = \tau_1 \tau_0 \\
    \rho_0 = \tau_\infty \tau_1 \\
    \rho_1 = \tau_0 \tau_\infty
  \end{cases}
  \qquad \rho_\infty \rho_1 \rho_0 \;\text{\struck{$\rho_\infty$}}\; = 1
\end{equation*}
]

\page{114}
\note{p.~2 de l'auteur.}

On introduit des sommets
\begin{equation*}
  \tag{6} R^{i,j}_r \qquad r \in \mathcal{R},\ (i,j) \in \mathrm{Rep}\{0,1,\infty\}
  \quad \text{i.e.\ } i,j \in \{0,1,\infty\},\ i \neq j
\end{equation*}
comme indiqué, satisfaisant aux seules relations
\begin{equation*}
  \tag{7} \boxed{R^{i,j}_r = R^{i,j}_{\tau_k(r)}, \quad k \notin \{i,j\}}
\end{equation*}
On dit que $R^{ij}_r$ est « proche » \add{du sommet} \add{$R^i_r$ de $r$} ($= R_s, R_a, R_f$ suivant le cas), il se trouve sur le côté \struck{\ill{}} \add{qui joint} \ldots{} \struck{\ill{}} au sommet de type $j$.

On introduit des arcs
\begin{equation*}
  \tag{8} \lambda^i_r : R^{i,i+1}_r \longrightarrow R^{i,i+2}_r \qquad r \in \mathcal{R},\ i \in \{0,1,\infty\}
\end{equation*}
\begin{equation*}
  \tag{9} a^i_r : R^{i+1,i}_r \;\text{\struck{\ill{}}}\; \longrightarrow R^{i+1,i+2}_r \qquad r \in \mathcal{R},\ i \in \{0,1,\infty\}
\end{equation*}
\struck{qui} les $a^i_r$ entourent la droite \uncertain{du} \uncertain{milieu} repère $r$ en formant une contour fermé dans le sens \struck{\ill{}} \add{défini par} l'orientation de $r$ \add{les $=$ désignent} \struck{\ill{}} les $\lambda^i_r$ \uncertain{sont} \uncertain{des} \uncertain{portions} \uncertain{d'arcs} \uncertain{autour} des sommets de $r$ du type $i$ dans le sens correspondant à l'orientation de $r$. (Cela \struck{\ill{}} signifie que : la convention de notation ne dépend que des sommets \struck{identifiés} \ill{} \ill{} \ill{}). On a les
\begin{equation*}
  \tag{10} \boxed{a^i_r = a^i_{\tau_i(r)}}
\end{equation*}

\marginal{\textbf{NB} On écrit \uncertain{pour} \uncertain{simplifier} \uncertain{(8)} \uncertain{et} \uncertain{(9)} \ldots{} \uncertain{(11)} \ldots{} en notant $i+1$, $i+2$ \uncertain{ou} \uncertain{plutôt} \uncertain{les} \uncertain{successeurs} de $i$ dans l'ordre circulaire $(0, 1, \infty)$, $(1, \infty, 0)$, $(\infty, 0, 1)$ \uncertain{identifiés} par $(i, i+1, i+2)$.}
\marginal{\textbf{NB} On \struck{désigne} \uncertain{appelle} \uncertain{flèches} les \ill{} \ill{} \uncertain{de} \ill{} \uncertain{sommets} \uncertain{et} \uncertain{arcs} \ldots{} \uncertain{les} \uncertain{trous} \uncertain{bordés} \uncertain{par} \uncertain{les} \uncertain{$R^i_r$} (\uncertain{trous} de \uncertain{dimension} $= 1$)}
\note{deux notes marginales écrites en diagonale dans la marge gauche ; lecture très partielle.}

\textbf{NB} \struck{On désigne} On \uncertain{observe} que la carte combinatoire à trous \uncertain{réalisée} par sa réalisation \uncertain{conforme} \uncertain{canonique}, qui est une surface $\mathcal{C}$ \uncertain{analytique} réelle dans $\mathbb{C}^\ast$, \uncertain{dont} on peut considérer \uncertain{la} \uncertain{réalisation} \uncertain{en} les $R_\varphi$ ($\varphi \in S$) et la \uncertain{découpe} de l'\uncertain{identité}, qui est une surface à bord analytique réelle (mais pas conforme) dont l'intérieur \uncertain{coïncide} avec $X \setminus S$, on la note
\note{la phrase se poursuit à la page suivante.}

\page{115}
$\tilde X^{\ast\ast}$ ; (les trois $\ast$ signifient qu'on a enlevé \uncertain{trois} \uncertain{types} \add{\uncertain{de} \uncertain{trous}} de points de type $0, 1, \infty$, des points \uncertain{comme} \uncertain{seuls} \add{$\infty$, $1$}, on définirait de même $\tilde X^{\cdot\ast}$, $\tilde X^{\ast\ast}$, $\tilde X^{\ast\cdot}$, etc. \struck{\ill{}} L'intérieur de cette surface à bord, savoir $X \setminus S = X \setminus \{R_\varphi \mid \varphi \in S\}$, est noté $X^{\ast\ast}$, on définit de même $X^{\ast\ast}$ etc. Les points $R^{ij}_r$ et les chemins précédents sont \uncertain{regardés}, la plus commodément, comme un point et chemins (\uncertain{parfaitement} canoniques \uncertain{pas} \uncertain{seulement} \uncertain{des} classes d'homotopie !) dans $\tilde X^{\ast\ast}$. On \uncertain{notera} par les $R_\varphi$ ($\varphi \in S$) un \uncertain{seul} \uncertain{point} des $\tilde X^{\ast\ast}_{\varphi}$
\uncertain{mais} dans l'application $C^\infty$ \uncertain{canoniques} \uncertain{analytiques} réelle) canonique
\begin{equation*}
  \tag{11} \tilde X^{\ast\ast} \longrightarrow X
\end{equation*}
\struck{On} \uncertain{pose} $R^{ij}_r$ \uncertain{va} \uncertain{sur} $R^i_r = R_\varphi$, où $\varphi$ est le \struck{sommet} \add{\uncertain{point}} de type $i$ de $r$.

\drawing{En marge gauche, un triangle de sommets $0$, $1$, $\infty$ ; à chaque sommet un petit arc fléché, et le long de la base une flèche intérieure marquée $b^0_r$.}

On a, dans $\Pi_1 X^{\ast\ast}$, la relation
\begin{equation*}
  \tag{12}
  \boxed{\underbrace{(\lambda^0_r)^{-1} a^1_r}_{\overset{\text{def}}{=}\, b^\infty_r}\;
  \underbrace{(\lambda^\infty_r)^{-1} a^0_r}_{b^1_r}\;
  \underbrace{(\lambda^1_r)^{-1} a^\infty_r}_{b^0_r} = 1}
  \qquad \text{lacet en } R^{1\infty}_r
\end{equation*}
Les sommets $R^{ij}_r$, les flèches $\lambda^i_r$ et $a^i_r$, et les relations (10) (12), forment une présentation du groupoïde $\Pi_1(\tilde X^{\ast\ast}, (R^{ij}_r))$.
\note{« on la note » (fin de p.~114) s'achève sur $\tilde X^{\ast\ast}$ en tête de cette page. Les exposants $\ast$, $\cdot$ sont lus comme à la p.~111, sans certitude sur leur disposition.}

\page{116}
\note{p.~3 de l'auteur.}

Les notations introduites sont telles qu'elles gardent un sens pour toute carte triangulaire « partielle » du type $\{0, 1, \infty\}$, \uncertain{sans} que celle-\uncertain{ci} \struck{\ill{}} soit \struck{la} \uncertain{de} \ill{} \uncertain{nécessairement} la subdivision barycentrique d'une carte $X$, condition qui s'exprime par la \struck{condition} \add{dernière} relation
\[
  \text{\struck{\ill{}}} \qquad (\tau_0 \tau_\infty)^2 = 1 \quad \text{i.e.} \quad \rho_1^2 = 1
\]
dans (4) (qui \uncertain{s'énonce} \uncertain{donc} par la suite) et qui signifie aussi que les \struck{points} \add{sommets} de type $1$ sont d'ordre réduit $2$ (d'ordre ordinaire $4$), on aura que
\[
  \tau_0 \tau_\infty(r) = \tau_\infty \tau_0(r) \qquad \text{pour tout } r.
\]
Dans le cas général, $\rho_1$ est donc une op\supplied{ération}. qui \uncertain{joue} un peu plus un rôle particulier \uncertain{pour} \ill{} : c'est la « rotation d'un cran » de $r$ autour du sommet de type $1$, dans le \struck{\ill{}} sens de \uncertain{l'orientation} de $r$.

On introduit
\begin{equation*}
  \tag{13}
  \begin{cases}
    \Lambda^i_r = (\lambda^i_{\tau_{i+1}(r)})^{-1} \lambda^i_r : R^{i,i+1}_r \longrightarrow R^{i,i+1}_{\tau_{i+1}(r)} \\
    \Lambda'^i_r = \lambda^i_{\tau_{i-1}(r)} (\lambda^i_r)^{-1} : R^{i,i-1}_r \longrightarrow R^{i,i-1}_{\tau_{i-1}(r)}
  \end{cases}
\end{equation*}
« rotation d'un double cran » autour du sommet de type $i$ de $r$, à partir du sommet proche de type $(i, i+1)$ \add{resp.\ $(i, i-1)$}, dans le sens de rotation donné \add{resp.\ au sens opposé} par l'orientation de $r$. On a les relations
\begin{equation*}
  \tag{14}
  \boxed{\Lambda^i_r \Lambda^i_{\tau_{i+1}(r)} = 1} \qquad
  \boxed{\Lambda'^i_r \Lambda'^i_{\tau_{i-1}(r)} = 1}
\end{equation*}
\note{les indices $\tau_{i\pm1}$ de (13) et (14) sont lus avec hésitation : le signe en indice est mal formé.}

\drawing{En bas à gauche, un morceau de carte : le triangle $r$ entre les sommets $0$, $1$, $\infty$ (deux exemplaires de $0$, deux de $1$, deux de $\infty$ sur les bords) ; sommets proches $R^{0,1}_r$, $R^{1,\infty}_r$, $R^{\infty,0}_r$, $R^{0}_{\tau_1 r}$, $R^{1\infty}_{\tau_0 r}$, $R^{\infty,0}_{\tau_0 r}$ ; arcs $\lambda^0_r$, $\lambda^1_r$, $\lambda^\infty_r$, $\lambda^0_{\tau_1 r}$, $\lambda^1_{\tau_0 r}$, $\lambda^\infty_{\tau_0 r}$ ; et les doubles crans $\Lambda^0_r$, $\Lambda^1_r$, $\Lambda^\infty_r$ en arcs courbes autour des sommets.}

\page{117}
\note{p.~4 de l'auteur ; en tête, « Suite ».}

\struck{L'introduction des $\Lambda^i_r$}

Fixons \uncertain{maintenant} un indice de $\{0, 1, \infty\}$, disons $1$ (\uncertain{conséquence} \uncertain{au} cas où la triangulation est barycentrique, \add{où} \uncertain{donc} $1$ joue un rôle particulier !). Nous pouvons éliminer les sommets $R^{ij}_r$ qui se trouvent sur des côtés de $r$ qui sont de type $1$, i.e.\ pour lesquels $i, j \neq 1$, et l'arc $a^1_r$ qui les joint, en gardant les sommets
\begin{equation*}
  \tag{14} \{R^{ij}_r\} \quad
  \begin{cases}
    r \in \mathcal{R} \\
    i \text{ ou } j \text{ est égal à } 1 \text{ i.e.\ } \{i,j\} \neq \{0,\infty\}
  \end{cases}
\end{equation*}
\note{il numérote (14) à nouveau, alors que (14) désigne déjà les relations de la p.~116.}
et les arcs
\begin{equation*}
  \tag{15}
  \begin{cases}
    \lambda^1_r & r \in \mathcal{R} \\
    \Lambda^0_r,\ \Lambda^\infty_r & r \in \mathcal{R},\ \text{\struck{\ill{}}} \\
    a^i_r & r \in \mathcal{R},\ i \in \{0, \infty\}
  \end{cases}
  \qquad \Lambda^0_r \Lambda^0_{\tau_1 r} =
\end{equation*}
avec comme relations les relations (14) de redondance \add{pour $i = 0$ et $i = \infty$} \struck{\uncertain{toutes} les relations du type $\Lambda^0_r = 1$} et

\begin{equation*}
  \tag{16}
  \text{\struck{$\Lambda^\infty_{\tau_1 r} (a_{\tau_1 r})^{-1} \lambda^1_{\tau_1 r} (a^\infty_{\tau_1 r})^{-1} \Lambda^0_r a^\infty_r \lambda^1_r a^0_r = 1$}}
\end{equation*}
\struck{pour $r \in \mathcal{R}$ (lacet en $R^{\infty 1}_r$)}
\struck{\textbf{NB} Cette relation est équivalente à la relation où $r$ est remplacé par $\tau_1 r$}
\note{tout ce bloc est barré de traits obliques ; sous le premier facteur, l'accolade $\beta^{1\,-1}_{\tau_1 r}$, sous le dernier $\beta^1_r$.}

\begin{equation*}
  \tag{16}
  \boxed{\Lambda'^\infty_{\tau_1(r)}
  \underbrace{\bigl(a^0_{\tau_1 r} (\lambda^1_{\tau_1 r})^{-1} a^\infty_{\tau_1 r}\bigr)}_{\beta^1_{\tau_1(r)}}{}^{\!-1}
  \Lambda^0_r
  \underbrace{\bigl((a^\infty_r)^{-1} \lambda^1_r (a^0_r)^{-1}\bigr)}_{\beta^{1\,-1}_r} = 1}
\end{equation*}
\marginal{relation stable par passage de $r$ à $\tau_1(r)$, donc il suffit de la prendre pour $r \in \mathcal{R}^+$ (si $X$ orientable)}
\note{l'exposant après la première parenthèse est surchargé ; lu $-1$ avec hésitation.}

\drawing{Au milieu à gauche, une étoile de triangles hachurés autour d'un sommet $1$, avec des sommets $0$, $\infty$ alternés sur le pourtour, chemins pointillés et flèches courbes de rotation ; les triangles $r$ et $\tau_1 r$ sont marqués.}
\marginal{\uncertain{rouge} par des \uncertain{carrés} ?, \uncertain{dont} \uncertain{seulement} \uncertain{seuls} \uncertain{les} \uncertain{sommets} de type $1$, \uncertain{sont} « arrondis », de façon à faire des \uncertain{losanges} \uncertain{évitant} les trous}

\drawing{En bas à gauche, un triangle $0$, $1$, $\infty$ et, à l'intérieur, le chemin $\beta^1_r$ fait de trois flèches.}

\page{118}
\drawing{Au crayon, un morceau de carte (cinq faces polygonales) ; le long d'une chaîne d'arêtes, de petits cercles fléchés autour des sommets, reliés par des flèches. Légende : « Cas IV ».}

\page{119}
\note{p.~5 de l'auteur.}

On a donc trouvé $4 = 1 + 3$ présentations du groupoïde fondamental $\Pi_1(\tilde X^{\ast\ast}, S)$, suivant qu'on utilise tous les chemins (8), (9), ou qu'on élimine les chemins $a^i_r$ \add{($i$ fixé, dans l'exemple $i = 1$)} et les $2$ sommets qui se trouvent dessus. Je vais maintenant examiner le cas où on élimine \struck{\ill{}} deux côtés des triangles fondamentaux, i.e.\ où on ne garde les $a^i_r$ que d'un seul type -- disons les $a^\infty_r$ -- ce qui conduit aussi : éliminons tous les sommets \add{et chemins qui \uncertain{vivent} \uncertain{sur} \uncertain{les} \uncertain{arêtes}} proches des $R_\varphi$ ($\varphi \in S$) de type \struck{$\infty$ \ill{}}, \struck{\ill{}, des contours des faces} \add{\uncertain{dans}} \struck{faces} -- i.e.\ \struck{les $R^i_r$ $j \neq \infty$} à ne garder que les $R^{ij}_r$ \add{$(i,j) \in \{0,1\}$} \struck{\ill{} pour avoir $j = \infty$} :
\begin{equation*}
  \tag{17} R^{ij}_r \quad
  \begin{cases}
    r \in \mathcal{R} \\
    \{i, j\} = \{0, 1\}
  \end{cases}
  \qquad\quad
  \boxed{
  \begin{cases}
    R^{0,1}_r = R^{0,1}_{\tau_\infty r} \\
    R^{1,0}_r = R^{1,0}_{\tau_\infty(r)}
  \end{cases}}
\end{equation*}
\marginal{relations de redondance \uncertain{pour} \uncertain{mémoire}}
\begin{equation*}
  \tag{18}
  \begin{cases}
    a^\infty_r : R^{0,1}_r \;\text{\struck{\ill{}}}\; \longrightarrow R^{1,0}_r & (r \in \mathcal{R}) \\
    \Lambda^0_r \;\text{\struck{$\Lambda^1_r$}} : R^{01}_r \longrightarrow R^{01}_{\tau_1(r)} \\
    \Lambda'^1_r : R^{10}_r \longrightarrow R^{1,0}_{\tau_0(r)}
  \end{cases}
\end{equation*}
\add{avec} relations de redondance seulement
\begin{equation*}
  \tag{19}
  \boxed{
  \begin{aligned}
    \Lambda^0_r \Lambda^0_{\tau_1(r)} &= 1 \\
    \Lambda'^1_r \Lambda'^1_{\tau_0(r)} &= 1
  \end{aligned}}
\end{equation*}
c'est donc directement exprimé comme un groupoïde libre. Cela fait trois autres présentations de $\Pi_1(X)$, mais sur des ens\supplied{embles} de sommets différents, savoir $\Pi_1(\tilde X^{\ast\ast}, S \setminus S_i)$, pour $i = \infty$ (cas explicité) ou $0$ ou $1$.

\drawing{En bas à gauche, une étoile de faces autour d'un sommet $\infty$, entourée d'autres sommets $0$, $1$, $\infty$ ; les arcs $\Lambda^0_r$, $\Lambda'^1_r$ en flèches courbes, chemins hachurés sur certaines arêtes.}
\marginal{\uncertain{rouge} par des \uncertain{faces} polygonales du centre des sommets de type $\infty$ -- les coins de ces polygones étant « arrondis » pour éviter les trous}

\page{120}
Je vais donner \struck{\ill{}} \add{trois} autres présentations, où on élimine non seulement les p\supplied{oin}ts proches (\uncertain{distincts}) des $R_\varphi$ sommets de type $\infty$, \add{et les sommets \uncertain{qui} \uncertain{leur} \uncertain{sont} \uncertain{liés} \uncertain{directement}} mais aussi ceux proches des sommets (distincts) de type $1$, en gardant donc que les sommets \add{$R^{ij}_r$} de type $0$, i.e.
\begin{equation*}
  \tag{20} R^{0,1}_r \qquad r \in \mathcal{R}
  \qquad\quad
  \boxed{R^{0,1}_r = R^{0,1}_{\tau_\infty r}}
\end{equation*}
\marginal{relations de redondance pour mémoire $(r \in \mathcal{R})$}
et on introduit les chemins
\begin{equation*}
  \tag{21}
  \begin{cases}
    A^{0,1}_r \overset{\text{def}}{=} (a^\infty_{\tau_0 r})^{-1} \;\text{\struck{\ill{}}}\; (\Lambda'^1_r)^{-1} a^\infty_r : R^{0,1}_r \longrightarrow R^{0,1}_{\tau_0(r)} \\
    \Lambda^0_r : R^{0,1}_r \longrightarrow R^{0,1}_{\tau_1 r} & (r \in \mathcal{R})
  \end{cases}
\end{equation*}
\marginal{plus les chemins \uncertain{déjà} \uncertain{connus}}
\note{la définition de $A^{0,1}_r$ est surchargée ; l'ordre et les exposants des facteurs sont lus avec hésitation.}
avec les relations de redondance
\begin{equation*}
  \tag{22}
  \boxed{
  \begin{cases}
    \Lambda^0_r \Lambda^0_{\tau_1(r)} = 1 & r \in \mathcal{R} \\
    A^{0,1}_r A^{0,1}_{\tau_0(r)} = 1 & r \in \mathcal{R}
  \end{cases}}
\end{equation*}
et c'est tout : $\Pi_1(\tilde X^{\ast\ast}, \{R^{01}_r\}_{r \in \mathcal{R}})$ est \add{présenté} \struck{engendré} par les sommets et flèches (20) (21), avec les seules relations de redondance (20) (22).

\drawing{En marge gauche, un triangle $r$ de sommets $0$, $1$, $\infty$, et son voisin $\tau_0 r$ ; le chemin $A^{0,1}_r$ va du point $R^{01}_r$ au point $R^{01}_{\tau_0 r}$ en passant autour du sommet $1$.}

\textbf{Bouchage de trous} Dans chacun des quatre cas envisagés, (en fait $10 = 1 + 3 + 3 + 3$ cas), le « bouchage » d'un trou $R_\varphi$ ($\varphi \in S = S_0 \sqcup S_1 \sqcup S_\infty$) se fait de façon essentiellement triviale. Dans chacun des quatre cas, il \uncertain{faut}, \uncertain{en} \uncertain{termes} des \struck{\ill{}} \uncertain{flèches} génératrices choisies, définir \struck{les}
\note{la page s'arrête au milieu de la phrase ; le développement sur le bouchage des trous se poursuit au-delà de ce lot.}

\end{document}
